\section{Modified Dispersion from Spectral Coefficients}
  \label{sec:dispersion}

  \subsection{Quadratic Operator and Pole Structure}
    \label{subsec:quadratic-operator-and-pole-structure}

    In the transverse-traceless sector established in Section~\ref{sec:ir_spectrum},
    the quadratic operator takes the factorized form
    \[
      \mathcal{O}_{\mathrm{TT}} = c_{\mathrm{EH}}\Box + \beta \Box^2 = c_{\mathrm{EH}}
      \Box
      \left( 1 + \frac{\beta}{c_{\mathrm{EH}}}\Box \right).
    \]

    The coefficient $\beta$ originates from the Weyl-squared term in the renormalized spectral expansion.
    For a Laplace-type scalar minimal operator in four dimensions, the Seeley--DeWitt coefficient $a_4$ contains
    \[
      \alpha_C = \frac{1}{360(4\pi)^2},
    \]
    multiplying $C_{\mu\nu\rho\sigma}^2$.
    Using
    \[
      \delta \int \sqrt g\, C^2 = -2 \int \sqrt g\, B_{\mu\nu}\,\delta g^{\mu\nu},
    \]
    one obtains
    \[
      \beta = -2\alpha_C = -\frac{1}{180(4\pi)^2}.
    \]
    The heat-kernel expansion and the associated Seeley--DeWitt
    coefficients are reviewed in~\cite{Vassilevich2003,DeWitt1965}.
    The variation of the Weyl-squared term leading to the Bach tensor is
    standard in higher-derivative gravity~\cite{Bach1921,Stelle1977}.

    The Einstein prefactor is parameterized as
    \[
      c_{\mathrm{EH}} = \zeta\,(4\pi)^{-2}\,\ell_\chi^{-2},
    \]
    where $\zeta=\mathcal{O}(1)$ accounts for scheme-dependent
    normalizations and possible multiplicity factors.
    In the natural spectral identification of the renormalization scale
    with the cutoff, $\mu=\ell_\chi^{-1}$, and using $a_2=R/6$, one has $c_{\mathrm{EH}}=\mu^2/[6(4\pi)^2]$,
    corresponding to $\zeta=1/6$.

  \subsection{Infrared Dispersion Relation}
    \label{subsec:infrared-dispersion-relation}

    For plane-wave solutions $h^{\mathrm{TT}}_{\mu\nu} \sim e^{ik\cdot x}$, the operator yields the dispersion equation
    \[
      -c_{\mathrm{EH}} k^2 + \beta k^4 = 0.
    \]

    Solving for $\omega^2$ gives
    \[
      \omega^2 = c^2 k^2 - \frac{\beta}{c_{\mathrm{EH}}} k^4.
    \]

    Using the expressions above,
    \[
      \frac{\beta}{c_{\mathrm{EH}}} = -\frac{1}{180\,\zeta}\, \ell_\chi^2,
    \]
    so that the infrared dispersion relation becomes
    \[
      \omega^2 = c^2 k^2 - \gamma \ell_\chi^2 k^4 + \mathcal{O}(k^6),
    \]
    with
    \[
      \gamma = \frac{1}{180\,\zeta}.
    \]

    In the natural spectral scheme $\zeta=1/6$,
    \[
      \gamma = \frac{1}{30},
    \]
    and therefore
    \[
      \omega^2 = c^2 k^2 - \frac{1}{30}\, \ell_\chi^2 k^4 + \mathcal{O}(k^6).
    \]

  \subsection{Physical Interpretation}
    \label{subsec:physical-interpretation}

    The correction term is negative, implying a slight sub-luminal deviation of the group velocity at high frequency,
    \[
      v_g = \frac{\partial \omega}{\partial k} \simeq c \left( 1 - \frac{3}{2}\gamma \ell_\chi^2 k^2 \right).
    \]

    Since the modification is suppressed by $(k\ell_\chi)^2$, it is negligible in the regime $k\ell_\chi \ll 1$.
    The additional pole at $k^2 \sim \ell_\chi^{-2}$ corresponds to a heavy excitation beyond the
    infrared window probed by current interferometers.

    The modified dispersion relation is therefore fully consistent with the infrared Einstein limit,
    while providing a concrete $k^4$ correction that can be directly confronted with gravitational-wave observations.
